\chapter{投资结论与行动框架}

\begin{houseviewbox}[核心判断｜盈利预告的广度不等于全面风险偏好]
我们不把1,680家公司的预告扩散解读为全面风险偏好开启。可执行机会集中在三类：第一，已经形成利润—现金流—外部锚闭环的可复核验证模型；第二，利润修复可信、但仍需下一季证据确认的事件验证标的；第三，价格已先行或分母质量不足、只能作为风险监控的标的。因此，本报告的重点不是寻找最高目标价，也不是直接给出配置建议，而是识别在当前价格下可被证据、外部锚和风险预算共同约束的验证队列。
\end{houseviewbox}

\section{投资行动摘要}

\noindent
\begin{exhibitbox}[图表 1｜当前可复核模型与行动含义]
\small
\begin{tabularx}{\textwidth}{L{1.1cm}L{1.7cm}L{1.5cm}R{1.0cm}R{1.0cm}R{1.1cm}X}
\toprule
\textbf{代码} & \textbf{标的} & \textbf{角色} & \textbf{现价} & \textbf{概率值} & \textbf{空间} & \textbf{行动与约束} \\
\midrule
000703 & 恒逸石化 & 可复核正向验证模型 & 14.48 & 23.42 & 61.7\% & 仅作验证模型；等待H2扣非、现金流和外部锚同步确认后再讨论风险暴露 \\
002379 & 宏桥控股 & 可复核正向验证模型 & 18.66 & 25.80 & 38.3\% & 仅作验证模型；等待H2扣非、现金流和外部锚同步确认后再讨论风险暴露 \\
300014 & 亿纬锂能 & 可复核正向验证模型 & 52.28 & 60.43 & 15.6\% & 仅作验证模型；等待H2扣非、现金流和外部锚同步确认后再讨论风险暴露 \\
000155 & 川能动力 & 下行纪律对照 & 11.09 & 10.04 & -9.5\% & 概率值低于现价；不新增风险暴露，用于约束追价和估值纪律 \\
\bottomrule
\end{tabularx}
\normalsize
\sourcenote{官方业绩预告、Q1财务、2026年7月15日收盘价及已归档原始研报；概率值为本机构情景值，非收益承诺。}
\end{exhibitbox}

\noindent
\begin{exhibitbox}[图表 2｜三层行动框架：从可复核价值到风险纪律]
\begin{minipage}[t]{0.31\textwidth}
\textbf{\color{riskgreen}可复核正向验证}
\par\smallskip
当前价格、情景分母、原始研报锚与风险触发器均可复核。
\par\smallskip
\textbf{代表：}恒逸石化、宏桥控股。
\par\smallskip
\textbf{纪律：}不直接配置；H2扣非或现金流偏离即降级。
\end{minipage}\hfill
\begin{minipage}[t]{0.31\textwidth}
\textbf{\color{riskamber}事件验证观察}
\par\smallskip
盈利修复有线索，但外部锚、现金流或持续性仍未闭环。
\par\smallskip
\textbf{代表：}中国船舶、中国铝业、三六零、中远海能。
\par\smallskip
\textbf{纪律：}不以模型空间作为行动理由。
\end{minipage}\hfill
\begin{minipage}[t]{0.31\textwidth}
\textbf{\color{riskred}估值纪律观察}
\par\smallskip
概率值低于现价、周期盈利可能回落或证据时效不足。
\par\smallskip
\textbf{代表：}川能动力及高估值样本。
\par\smallskip
\textbf{纪律：}不追价，作为行业和风险偏好的反向指标。
\end{minipage}
\sourcenote{行动分层由当前价格、估值分母、外部锚和证据质量共同决定，不等同于券商投资评级。}
\end{exhibitbox}

\section{执行纪律：先验证，再放大}

可复核正向模型只在利润、现金流与外部锚三项同时未偏离时维持；任一项失效即转入事件验证。事件验证观察不以模型空间作为建仓理由，须等待下一份财报、订单或价格信号把假设变为事实。估值纪律观察不追价，主要用于识别行业盈利、价格与风险偏好是否发生反向变化。

\section{补证后的候选池调整}

原中等证据标的共48只，补证后基础证据不再是主要阻断项；其中6只进入优先验证候选、11只进入扩展事件验证、3只进入安全边际观察，28只仅保留为风险监控或暂不纳入。没有任何原中等证据标的直接进入正式模型池，核心原因仍是当前可正权外部估值锚不足或价格纪律不满足。

\noindent
\begin{exhibitbox}[图表 3｜补证后候选准入摘要]
\scriptsize
\begin{tabularx}{\textwidth}{L{2.5cm}R{0.9cm}L{4.0cm}X}
\toprule
\textbf{层级} & \textbf{数量} & \textbf{代表标的} & \textbf{使用方式} \\
\midrule
优先验证候选 & 6 & 吉林敖东、辽宁成大、中信金属、三六零、宁波华翔、海德股份 & 进入优先验证队列；等待外部锚、H2扣非和现金流确认 \\
扩展事件验证 & 11 & 九安医疗、焦作万方、神火股份、香农芯创、海正药业、高德红外 & 列入扩展观察，不因空间大直接升级 \\
安全边际观察 & 3 & 中远海能、赣锋锂业、越秀资本 & 低优先级跟踪，等待安全边际扩大 \\
风险/不纳入 & 28 & 渤海租赁、翔鹭钨业、申通快递、长江证券、中泰证券、上海石化 & 概率值低于现价或现金流/外部锚不足，不新增为候选 \\
\bottomrule
\end{tabularx}
\normalsize
\sourcenote{完整逐票准入结论见附录C；本表为投委会摘要，不构成正式推荐。}
\end{exhibitbox}

\section{为什么没有按100\%以上空间直接入选}

全142只候选中，本轮模型空间超过100\%的标的共有6只，其中3只进入39只优先池，0只进入正式模型。专项核验已对6/6只完成原始API、原始PDF、目标价字段与现金流闭环；当前正权原始目标为0只。因此未升级不是资料空白，而是已核验证据不支持把内部高空间转换为正式模型。

\noindent
\begin{exhibitbox}[图表 3A｜100\%以上空间样本的拦截链条]
\scriptsize
\begin{tabularx}{\textwidth}{L{1.05cm}L{1.55cm}R{1.0cm}L{1.35cm}L{1.35cm}X}
\toprule
\textbf{代码} & \textbf{标的} & \textbf{空间} & \textbf{优先池} & \textbf{正式池} & \textbf{未直接入选原因} \\
\midrule
002432 & 九安医疗 & 323.4\% & 否 & 否 & 未入优先池；当前原始目标0；Q1现金流-0.92亿元；保留候选观察，维持未进入优先池及非正式模型 \\
000623 & 吉林敖东 & 317.0\% & 是 & 否 & 优先池内；当前原始目标0；Q1现金流0.11亿元；保留优先池验证，维持非正式模型 \\
600739 & 辽宁成大 & 300.0\% & 是 & 否 & 优先池内；当前原始目标0；Q1现金流-1.76亿元；保留优先池验证，维持非正式模型 \\
000685 & 中山公用 & 162.8\% & 否 & 否 & 未入优先池；当前原始目标0；Q1现金流-0.27亿元；保留候选观察，维持未进入优先池及非正式模型 \\
600150 & 中国船舶 & 142.9\% & 是 & 否 & 优先池内；当前原始目标0；Q1现金流81.52亿元；保留优先池验证，维持非正式模型 \\
301308 & 江波龙 & 119.2\% & 否 & 否 & 未入优先池；当前原始目标0；Q1现金流-28.75亿元；保留候选观察，维持未进入优先池及非正式模型 \\
\bottomrule
\end{tabularx}
\normalsize
\sourcenote{专项包已逐票核验东方财富原始API、已归档原始PDF、目标价字段与Q1现金流。媒体转载、聚合页和失败探针权重均为0。}
\end{exhibitbox}


\section{重点事件验证清单}

下表不是推荐名单，而是最需要在下一轮信息披露中验证的高敏感样本。其共同点是：模型空间可能存在，但当前尚不满足正式审计模型条件。

\noindent
\begin{exhibitbox}[图表 3B｜条件性机会：从模型空间到证据升级]
\scriptsize
\begin{tabularx}{\textwidth}{L{1.1cm}L{1.7cm}L{1.6cm}R{1.0cm}R{1.0cm}X}
\toprule
\textbf{代码} & \textbf{标的} & \textbf{估值框架} & \textbf{概率值} & \textbf{空间} & \textbf{必须验证的事实} \\
\midrule
600150 & 中国船舶 & 成长盈利情景 PE & 83.33 & 142.9\% & 订单交付、合同负债、产品结构和应收账款回款；反证：订单确认延后、交付节奏放缓、应收账款继续扩张或估值透支交付 \\
601600 & 中国铝业 & 周期调整扣非 EPS / PE & 14.07 & 56.7\% & 金属价格、产量、单位成本、库存和经营现金流；反证：价格回落、成本曲线抬升、产量不达预期或利润无法转化为现金 \\
601360 & 三六零 & 2026E 收入 PS & 13.10 & 45.6\% & 合同订单、项目交付、回款、研发投入和利润率；反证：订单确认延后、回款周期拉长、研发费用率上升或估值继续压缩 \\
600026 & 中远海能 & 周期调整扣非 EPS / PE & 17.45 & 10.2\% & 运价、运力利用率、燃料成本、资本开支和现金流；反证：运价回落、利用率下降、成本上升或周期盈利被错误年化 \\
002048 & 宁波华翔 & 正常化 EPS / PE & 28.50 & 33.0\% & 销量、单车盈利、产品结构、价格竞争和应收回款；反证：销量或单车盈利下滑、价格战加剧、产品结构恶化或现金流转弱 \\
002812 & 恩捷股份 & 产能周期正常化前瞻 EPS / PE & 85.29 & 55.2\% & 出货量、订单、产品价格、产能利用率和营运资金；反证：订单转化慢、价格竞争、利用率不足或库存/应收继续占用现金 \\
300390 & 天华新能 & 周期调整 EPS / PE & 77.30 & 28.8\% & 出货量、订单、产品价格、产能利用率和营运资金；反证：订单转化慢、价格竞争、利用率不足或库存/应收继续占用现金 \\
002240 & 盛新锂能 & 周期调整 EPS / PE & 37.89 & 21.8\% & 金属价格、产量、单位成本、库存和经营现金流；反证：价格回落、成本曲线抬升、产量不达预期或利润无法转化为现金 \\
\bottomrule
\end{tabularx}
\normalsize
\sourcenote{条件性模型仅用于安排验证优先级。若缺乏当前外部锚或现金流确认，则不升级为可复核正向验证。}
\end{exhibitbox}

本报告的反方情景同样明确：若H2扣非利润不延续、经营现金流继续背离、商品价格或行业价差回落，当前候选池中的高空间模型将被迅速下调，而不是通过提高倍数维持结论。
